\documentclass[a4paper,11pt]{article}
\usepackage{exam-style-842-unit}
\examinehead{842 信息化概论 · 第 8 单元练习册}{练习一 参考答案与解析}
\begin{document}

\unitanswerscover{842 信息化概论 · 单元练习册}{第 8 单元参考答案}{练习一（电子商务基础与模式）的答案与解析；每道判断题均给出改正后的正确表述。教材页码按“扫描页−7＝印刷页”换算，引文以 OCR 原文为准。}

\answerpart{一、单项选择题（10 题，每题 3 分，共 30 分）}

\q A。\begin{analysis}1997 年 1 月 6—7 日在法国首都巴黎，国际商会举行了世界电子商务会议（The World Business Agenda for Electronic Commerce），全世界商业、信息技术、法律等领域的专家和政府部门的代表共同讨论电子商务概念，这是目前电子商务较为权威的概念表述（教材 8.1.1，第 122 页）。B 世界贸易组织在电子商务专题报告中给出定义；C 联合国经济合作和发展组织（OECD）是较早系统研究电子商务的机构；D 全球信息基础设施委员会（GIIC）电子商务工作委员会提出报告草案，三者都不是“巴黎世界电子商务会议”。\end{analysis}

\q B。\begin{analysis}广义上讲，电子商务（Electronic Business，E-Business）就是通过电子手段进行的商业事务活动（教材 8.1.1，第 123 页）。A 是狭义电子商务（E-Commerce）的含义；C、D 都只截取了电子商务的某一部分，范围过窄。\end{analysis}

\q A。\begin{analysis}狭义上讲，电子商务（Electronic Commerce，E-Commerce）是指通过使用互联网完成的商务交易活动，即单指在网络上做买卖（教材 8.1.1，第 123 页）。B 是广义电子商务的范围；C 企业内部管理、D 政府网上采购都不是狭义电子商务的定义。\end{analysis}

\q A。\begin{analysis}IBM 提出了一个电子商务的定义公式，即电子商务＝Web＋IT，强调在网络计算环境下的商业化应用，把买方、卖方、厂商及其合作伙伴在因特网、内联网和外联网结合起来（教材 8.1.1，第 123 页）。B 中的 EDI 是 GE 描述企业间电子商务时强调的核心技术；C 物流与资金流是电子商务整合的对象；D 不是教材中的公式。\end{analysis}

\q A。\begin{analysis}B2B 方式是电子商务应用最多和最受企业重视的形式，其代表是马云的阿里巴巴电子商务模式（教材 8.1.1，第 127 页）。B 淘宝、C eBay 是 C2C 的代表；D 京东是典型的 B2C 平台。\end{analysis}

\q B。\begin{analysis}C2C 是消费者与消费者之间的电子商务，C2C 商务平台通过为买卖双方提供在线交易平台，使卖方可以主动提供商品上网拍卖，买方可以自行选择商品进行竞价，其代表是 eBay、淘宝电子商务模式（教材 8.1.1，第 128 页）。A 阿里巴巴是 B2B 代表；C 京东、D 苏宁易购以 B2C 为主。\end{analysis}

\q B。\begin{analysis}O2O（Online to Offline）是新兴起的一种电子商务模式，即将线下商务的机会与互联网结合在一起，让互联网成为线下交易的前台，线下服务用线上来揽客、消费者用线上筛选服务并在线结算，其最重要的特点是推广效果可查、每笔交易可跟踪（教材 8.1.1，第 130 页）。A 是 B2B；C 教材未作为独立模式列出；D 是 B2G。\end{analysis}

\q D。\begin{analysis}与传统商务相比，电子商务具有业务全球化、服务个性化、业务集成性、机会的均等性 4 个特点（教材 8.1.1，第 124—125 页）。“交易垄断性”与教材表述相反：互联网是开放的大市场，电子商务使中小企业无需庞大商业体系即可打开国际市场，机会是均等的。\end{analysis}

\q A。\begin{analysis}迈克尔·波特（Michael E. Porter）于 1985 年在其所著《竞争优势》一书中认为，任何企业的价值链都由一系列相互联系的创造价值的活动构成，可分为基本活动和辅助活动两类（教材 8.1.1，第 126 页）。B 布鲁克斯著《人月神话》；C 罗伊斯提出瀑布模型；D 梅棹忠夫首次提出“信息社会”概念。\end{analysis}

\q B。\begin{analysis}电子商务平台体系架构中，直接面向用户、负责页面展示与交互的是表示层（前端），通常包括 Web 网站、移动 App、小程序等多端接入（教材未设“电子商务平台体系架构”专节，按考试大纲“四、电子商务”补充）。A 数据层负责数据存储与管理；C 基础设施层提供网络、服务器、存储等资源；D 网络层负责数据传输，均不直接面向用户展示。\end{analysis}

\answerpart{二、多项选择题（4 题，每题 4 分，共 16 分）}

\q ABCD。\begin{analysis}电子商务的特点有 4 个：业务全球化、服务个性化、业务集成性、机会的均等性（教材 8.1.1，第 124—125 页）。业务全球化使交易不受时间空间制约；服务个性化使客户服务更方便；业务集成性表现为企业事务处理的整体性、与客户的互动性、与上下游伙伴的密切合作；机会的均等性使中小企业与大企业有同样多的机会。四项全对。\end{analysis}

\q ABCD。\begin{analysis}B2C 通常有 4 种经营模式：虚拟社群（Virtual Communities）、交易聚合（Transaction Aggregators）、广告网络（Advertising Network）、线上与线下结合的模式（O2O 模式）（教材 8.1.1，第 128 页）。四项全对。\end{analysis}

\q ABCD。\begin{analysis}教材在 B2B、B2C、C2C 之后指出，随着互联网应用不断拓展，一些新兴电子商务模式如 B2G、BOB、O2O 等涌现出来，并具体介绍了 B2G、C2B、BOB、B2T、O2O、ABC、O2P 等模式（教材 8.1.1，第 129—130 页）。A B2G 是企业与政府之间；B C2B 是消费者与企业之间；C O2O 是线上与线下结合；D B2T 是团队向商家采购（团购）。四项均属新兴模式。\end{analysis}

\q ABCD。\begin{analysis}电子商务平台体系架构通常分为：表示层（Web、App、小程序多端接入）、应用业务层（商品、订单、支付、物流、营销等业务功能）、数据层（数据库、数据仓库与数据分析）、基础设施层（网络、服务器、存储与云资源）；安全与运营管理体系贯穿各层（教材未设专节，按考试大纲“电子商务平台体系架构”补充）。四项全对。\end{analysis}

\answerpart{三、判断题（6 题，每题 2 分，共 12 分）}

\q ×。改正：电子商务有广义与狭义之分，网上购物只是狭义电子商务的典型形式；广义电子商务是通过电子手段进行的全部商业事务活动。\begin{analysis}教材 8.1.1 明确区分：广义电子商务（E-Business）是通过电子手段进行的商业事务活动，贯穿企业行为全过程；狭义电子商务（E-Commerce）仅指通过互联网完成的商务交易活动。把电子商务等同于网上购物，是以偏概全。\end{analysis}

\q √。\begin{analysis}教材 8.1.1 指出，B2B 方式是电子商务应用最多和最受企业重视的形式，企业可以使用 Internet 或其他网络对每笔交易寻找最佳合作伙伴，完成从订购到结算的全部交易行为（第 127 页）。\end{analysis}

\q ×。改正：C2C 平台经营者不负责物流，而是协助汇集市场资讯、建立信用评价制度等，物流由买卖双方自行解决。\begin{analysis}教材 8.1.1 在介绍 C2C 竞标网站时明确指出：网站经营者不负责物流，而是协助汇集市场资讯，以及建立信用评价等制度（第 128 页）。\end{analysis}

\q √。\begin{analysis}教材 8.1.1 指出：E-Business 包括的范围要比 E-Commerce 宽广得多，因为后者仅指简单的商务交易应用，即单指在网络上做买卖；而 E-Business 是存在于企业与企业之间、企业与客户之间、企业内部的一种联系网络，贯穿于企业行为的全过程（第 123 页）。\end{analysis}

\q √。\begin{analysis}教材 8.1.1 指出，互联网代表了一个开放性的大市场，使小型企业无需庞大的商业体系、昂贵的广告费用和众多的营销人员，只需网页就可以打开国际市场；电子商务的均等性对于中小型企业来说尤其有利（第 125 页）。\end{analysis}

\q ×。改正：价值链并不是孤立存在的，而是存在于由供应商价值链、企业价值链、渠道价值链和买方价值链共同构成的价值链系统中；企业与企业的竞争不只是某个环节的竞争，而是整个价值链的竞争。\begin{analysis}教材 8.1.1 在介绍波特价值链理论时明确指出价值链系统的构成，并强调整个价值链的综合竞争力决定企业的竞争力（第 126 页）。把价值链局限于企业内部，是对价值链理论的误读。\end{analysis}

\answerpart{四、配对题（6 对，每对 1 分，共 6 分）}

\q （1）—A\quad （2）—B\quad （3）—C\quad （4）—D\quad （5）—E\quad （6）—F。
\begin{analysis}B2B 代表为阿里巴巴；C2C 代表为 eBay、淘宝；O2O 是线上揽客、线下消费体验；1997 年世界电子商务会议在巴黎由国际商会举行；迈克尔·波特 1985 年在《竞争优势》中提出价值链理论；IBM 提出电子商务＝Web＋IT 的定义公式（教材 8.1.1）。\end{analysis}

\answerpart{五、简答题（4 题，每题 6 分，共 24 分）}

\q 参考答案要点：
（1）权威表述（1.5 分）：1997 年巴黎世界电子商务会议上，与会代表认为电子商务是指对整个贸易活动实现电子化；从涵盖范围看，是交易各方以电子交易方式而不是通过当面交换或直接面谈方式进行的任何形式的商业交易；从技术看，是一种多技术的集合体，包括交换数据（电子数据交换、电子邮件）、获得数据（共享数据库、电子公告牌）和自动捕获数据（条形码）等。
（2）广义与狭义（2 分）：广义电子商务（E-Business）是通过电子手段进行的商业事务活动，贯穿企业行为全过程，范围更宽；狭义电子商务（E-Commerce）是指通过使用互联网完成的商务交易活动，即单指在网络上做买卖。
（3）多方定义（1.5 分）：欧洲议会、美国政府《全球电子商务纲要》、OECD、WTO、GIIC 等分别从电子方式、Internet、电子化手段、电信网络、电子通信等角度给出定义；IBM 提出“电子商务＝Web＋IT”，HP 提出“电子化世界”。
（4）综合定义（1 分）：电子商务是依托现代信息技术和网络技术，集金融电子化、管理信息化、商贸网络化于一体，旨在实现物流、资金流与信息流和谐统一的新型贸易方式，以商务活动为主体，以计算机网络为基础，以电子化方式为手段，在法律许可范围内通过信息化手段进行商业事务处理和商务交易活动。（教材 8.1.1，第 122—123 页）

\q 参考答案要点：
电子商务将传统商务活动中的物流、资金流和信息流利用网络技术进行整合，与传统商务相比具有 4 个特点（每点 1.5 分）：
（1）业务全球化：交易各方可通过互动方式直接在网上完成交易及全部相关活动，商品和信息交换不再受时间、空间制约，企业可把商务活动范围扩展到全球，消费者的购物选择也是全球性的。
（2）服务个性化：客户不再受地域限制，服务质量成为商务活动关键；企业把客户服务过程移至万维网，客户可以足不出户、实时完成转账、查询、购物等操作，方便性突出。
（3）业务集成性：表现为企业事务处理的整体性、统一性，与客户的直接互动性，以及企业与销售方、供货方和商务伙伴间更密切的合作关系，从而规范流程、提高效率、降低成本。
（4）机会的均等性：互联网开放的大市场使中小企业无需庞大商业体系、昂贵广告费用和众多营销人员，只需网页即可打开国际市场，机会对大中小企业均等，对中小企业尤其有利。（教材 8.1.1，第 124—125 页）

\q 参考答案要点：
（1）B2B（2 分）：企业与企业之间的电子商务（Business to Business），以企业为主体，企业可以使用 Internet 或其他网络对每笔交易寻找最佳合作伙伴，完成从订购到结算的全部交易行为；是电子商务应用最多、最受企业重视的形式，也是企业改善竞争条件、建立竞争优势的主要方法，能带来更低价格、更高生产率和更低劳动成本；代表为阿里巴巴。
（2）B2C（2 分）：企业与消费者之间的电子商务（Business to Customer），企业通过网络销售产品或服务给个人消费者，类同于商业电子化的零售商务，如网络购物、证券公司网络下单、网站资料查询等；通常有虚拟社群、交易聚合、广告网络、线上与线下结合（O2O）4 种经营模式。
（3）C2C（2 分）：消费者与消费者之间的电子商务（Consumer to Consumer），C2C 商务平台为买卖双方提供在线交易平台，卖方主动提供商品上网拍卖，买方自行选择商品竞价；平台经营者不负责物流，而是协助汇集市场资讯、建立信用评价等制度；代表为 eBay、淘宝。
（教材 8.1.1，第 127—128 页；另见表 8-1 三种模式比较，第 128—129 页）

\q 参考答案要点（教材未设“电子商务平台体系架构”专节，按考试大纲“四、电子商务”补充；教材 8.1.1 的商业模式、价值链与模式分类可作基础）：
电子商务平台一般按分层架构设计（每层 1.5 分，共 6 分）：
（1）表示层（接入层）：面向用户提供 Web 网站、移动 App、微信小程序、第三方平台开放接口等多端入口，负责页面展示、交互与输入校验。
（2）应用业务层：实现商品管理、用户与会员、订单与交易、支付结算、营销推广、库存与物流、售后服务、数据分析等业务功能，是平台的核心。
（3）平台服务层：以组件或服务方式提供支付、物流、客服、消息、搜索、风控等可复用能力（教材中的网络银行、第三方支付、电商物流与供应链管理可视为该层外部服务）。
（4）数据层与基础设施层：数据库、数据仓库与大数据分析支撑经营决策；网络、服务器、存储与云资源提供运行基础。
（5）安全与运营管理体系：贯穿各层，包括身份认证、数据加密、支付安全、隐私保护、运维监控与平台规则治理。
分层的作用：职责分离、松耦合、便于扩展与维护，支撑教材所述物流、资金流、信息流的统一（8.1.1）。

\answerpart{六、综合题（1 题，共 12 分）}

\q 参考答案要点：
（1）模式判断（5 分，每项 1 分）：
企业自营旗舰店——企业通过网络向个人消费者销售产品，属 B2C；
合作社在平台开店——合作社作为企业性经营主体向消费者销售，属 B2C（若面向其他企业批量供货则属 B2B，需按实际交易对象判断）；
消费者之间的二手交易——属 C2C；
“线上下单、到果园体验采摘”——线上揽客、线下体验消费，属 O2O；
政府消费帮扶采购——政府向企业（合作社、农户组织）采购，属 B2G。
（2）价值链与价值网分析（4 分，每点 1 分，答出 4 点即可）：
第一，平台把分散的农户、合作社与消费者连接起来，缩短了传统多级批发环节，降低了交易成本。
第二，平台汇集市场资讯与信用评价，缓解信息不对称，使农产品能够以更合理的价格成交（对应教材 C2C/B2C 平台的作用）。
第三，按波特价值链理论，平台通过订单、支付、物流与售后服务等活动为顾客创造价值，并借助数据与营销开展辅助活动，提升整体竞争优势。
第四，电商使企业价值链呈现虚拟化、网络化特点，传统线性价值链转化为价值网：平台与供应商（合作社）、分销商（物流企业）、顾客共享信息、协同运作，快速响应市场需求。
第五，价值网以网络和信息系统为手段，使价值链上所有成员密切合作，快捷、可靠、高效地创造更多价值，最终让农户增收、合作社扩大销路、消费者获得可追溯的优质产品。
（3）平台体系架构（3 分）：可设计为“表示层—应用业务层—数据层—基础设施层＋安全运营体系”。表示层提供 Web、App、小程序与第三方平台接口；应用业务层实现商品、订单、支付、物流、营销、溯源等功能；数据层存储商品、订单、用户与溯源数据并支持分析；基础设施层提供网络、服务器与存储；安全与运营体系负责认证、支付安全、隐私保护与平台治理。该架构支撑物流、资金流、信息流的统一，符合大纲对电商平台体系架构的要求（教材未设专节，按大纲补充）。

\end{document}
